\section{Localization, frozen rows, and the side estimates}
\label{sec:localization}

Fix a later accuracy level \(t\ge12\), write \(k=k_t\), and choose a
final base segment \(R_0\in\mathcal R_{\rm base}\) and a pattern
\(z\in\mathcal P\) from Proposition~\ref{prop:hbase-data}.  With
\(\eta=e^{-k^\theta}\) and \(0<\theta<1\) as in
\eqref{eq:hbase-quality}, its target \eqref{eq:h-active-pattern-target} is
\[
 \mathfrak M_{R_0}(z)=\sum_{I\in\mathcal A_{R_0}}|I|
       \min\{\eta,|H_I(z_I)|\}.
\]
Here \(\mathcal A_{R_0}\) is the original active collection assigned
to this segment.  On a cell of depth \(s\), write
\(z_I=(z_{0,s},\ldots,z_{3,s})\) for its whole current arguments.
The truncation was obtained from the full tuple
\(Y\) before the \(B/U\) patterns were separated.  The pair estimate
for \(Y\) will be used after filter transfer restores the two complementary
\(B\) inputs to \(Y\).

We first work on one localization attempt with root \(R\subseteq R_0\).
Let \(\mathcal A_R^{\rm att}\subseteq\mathcal A_{R_0}\) be the original
active cells assigned to it.  Let \(\mathcal A_R^{\rm main}\) be
the cells in \(\mathcal A_R^{\rm att}\) surviving the initial
marking-failure cuts, namely those not at or below a marking-failure
root.  This collection is fixed before the dependent cuts.  Let
\(\mathcal A_R^{\rm ret}\subseteq\mathcal A_R^{\rm main}\) be those
retained after all cuts.  An assigned proper ancestor of a cut remains one
whole \(H_I\) summand.  The estimates in this section bound
the contribution of \(\mathcal A_R^{\rm ret}\) by a multiple of
\(|R|\).  A realized finite operator family is fixed during its
operator tests.  A later cut transfers original summands at that cell
and below to a new attempt without changing the parent's fixed
functions or operators.  Section~\ref{sec:completion} renews the
estimates on the new roots and sums their lengths.

All depths \(s,b,u_b,v_b,t_b\) below remain the original depths in
\(J\).  Passing to \(R\) does not reset an input schedule or its
auxiliary nearest-endpoint extension.  An endpoint output before
depth \(s(R)\) is set to zero.  The localization attempts here are
distinct from the fit attempts that produced the base segments.

\subsection{One attempt and its localization data}

Set \(P=\lceil\eta^{-\alpha}\rceil\), with the fixed \(\alpha>0\)
chosen in Section~\ref{sec:completion}.  For fixed \(\alpha\),
\(\log(2P)=O(k)\).  The main catalog consists of
the basic columns and the additional columns used by H's local
filters; its budget \(D_{\rm m}\) also anticipates the later helper
columns.  It bounds path counts, column sizes, coefficient bounds,
inverse private weights, and main-use multiplicities.  We use
\[
 \log D_{\rm base},\ \log D_{\rm m}\le C_qk^C,\qquad
 \mathcal L=k^{O(1)},\qquad
 G_0=\exp\!\bigl(C_q(1+\log(2k))^C\bigr).
\]
There are \(\mathcal L\) colors, and the relevant main histories have
\(O(\log\mathcal L)\) factors.  The value of \(G_0\) may increase
within this displayed class.

We describe the main data of \cite[Section~6]{H} before checking
their realization for the present inputs.  A present vertex is an
occurrence of an original basic column, with H's full normalized
curvature tag, including its horizontal chart data.  The present
lists grow by inheritance along downward paths.  In slot \(l\),
the quadratic chart phase is \(\kappa_l\) times the phase represented
by that tag, where \(\kappa=(-1,3,-3,1)\); linear modulation is not
part of the tag.  Distances are measured in H's ambient graph on all
hypothetical tags satisfying its bounds.  They depend on the tags and
absolute depth, and do not increase as depth increases.  Every added main
or helper column is attached to a present vertex and has its tag;
these insertions add no graph vertices.

A group \(G\) is a source index with a fixed color and a set of
actual central uses.  Its pad mask \(\mathcal M_{G,h}(I)\) is a set of
present vertices on the cell \(I\), retaining historical central
uses.  Such masks may continue by inheritance after the group has
ceased to be current on a branch.  At a given rank, the associated
pad projection is the orthogonal projection in \(L^2(I;\mathcal H)\) onto
constant-coefficient combinations of the available columns of at
most that rank whose tags belong to the mask.  Pads are nested on a
fixed cell, and each individual pad family is inherited on the
complete depth partitions and unused-cell continuations needed
below.  A positive smeared filter is the normalized average of its
\(P\) pad projections.  A named history is a specified ordered
history of these source filters and their complements.  After the
catalogs are fixed it is a fixed bounded operator; expanding its
complements and choosing individual pads gives ordered products of
inherited orthogonal projections.

Usability is relative to the pad mask supplying a range: that mask
is nonempty, and the cell has an actual use of \(G\) at or below it.
The \emph{Mask inclusion, closure, and separation} Lemma of
\cite[Section~6]{H} gives three properties.  At an actual use, the
pad-\(h\) mask contains every present vertex at graph distance at
most \(h\) from the center; its pad-\(h+1\) mask contains every
neighbor of a pad-\(h\) vertex.  For distinct groups of the same
color, vertices in their respective pad masks on overlapping usable
endpoints are separated by graph distance greater than one at the
finer endpoint.

We now verify the closure input.  At the cell \(I\) of depth \(s\)
being tested, hold fixed its one current basic argument
\(b=z_{l,s}\), of type \(B\) or \(U\).  For two specified positions,
let \(\mathfrak d_I(v,w)\) be the supremum of the absolute form with
those entries and arbitrary normalized \(L^2(I)\) functions in the
other positions.  At rank \(r\), H requires, for every named target
history and each of its prescribed individual-pad choices,
\begin{equation}
 \mathfrak d_I\bigl(\hbox{one lower-rank source},
       (1-P_{r,\mathcal M})\mathcal H b\bigr)
 \le \frac{\tau_3}{D_{\rm base}\sqrt{n_{r-1}}}\|b\|_{2,I},
 \qquad \tau_3=D_{\rm base}^{-C_\tau}.
 \label{eq:h-closure-test}
\end{equation}
The source is one pure available lower-rank column whose tag lies
in the target mask, and \(n_{r-1}\) bounds the path count through
the preceding rank.  This is Equation~\Heq{13} of
\cite[Section~6]{H}; the individual tests also control the
normalized averages.

In the \emph{Catalog closure and its budget} Lemma of H, the tested
state \(X_i=\mathcal H_i b\) changes with the insertion index \(i\)
because its projections acquire columns, while \(b\) stays fixed
on \(I\).  The \emph{Uniform same-tag test} Lemma supplies on failure
a fresh target-column increment satisfying
\[
 \|(P_{i,\rm new}-P_{i,\rm old})X_i\|_{2,I}
 \ge c\frac{\tau_3}{D_{\rm base}\sqrt{n_{r-1}}H_0}\|b\|_{2,I},
 \qquad H_0\le\exp(C_qk^C).
\]
For one named history and mask in a trial of at most \(T\)
insertions, the target increments are orthogonal rows and every
underlying positive projection is nested in \(i\).  Expanding
complements and applying the products-of-prefixes calculation of
\cite[Section~3]{H} gives
\begin{equation}
 \sum_i\|(P_{i,\rm new}-P_{i,\rm old})\mathcal H_i b\|_{2,I}^2
 \le (C\log(2T))^{O(\log\mathcal L)}\|b\|_{2,I}^2.
 \label{eq:h-closure-prefix}
\end{equation}
This is a fixed-vector estimate on \(b\), so it applies when \(b\)
contains \(F_{l,s}\): the depth \(s\) is fixed during the cell test.
The lower bound and \eqref{eq:h-closure-prefix} cancel the norm of
\(b\).  Counting named masks and histories and taking
\(n_r/n_{r-1}=\exp(k^{C_2})\) through \(\mathcal L+5\) ranks gives
\(\log n_{\max}\le C_qk^C\), using the active path count from the
base reduction.

This is an a priori count from names of contexts, including the
anticipated number of helper bases; it is independent of the
numeric edge height \(H_{\rm ed}\) and later mode precision.  As
in H, choose the mismatch accuracy after this count and then choose
\(H_{\rm ed}\) and realize compatible graph, mask, and main-catalog
data.  We fix these main projections, color histories, ball masks,
and actual central uses, including the small marking-failure cuts
and the inherited continuations with local refits on unused branches
needed for operator estimates.  The actual helper selections will
be made after their bases have been constructed.

The \emph{Bessel separation of usable endpoint spaces} Lemma of
\cite[Section~6]{H} supplies the following operator interface.  For
one fixed slot and color, let
\[
 h_G=\sum_{v,U}a_{v,G,U}\mathbf1_U\Phi_v
\]
be finite sums of available pure main columns.  Each coefficient is
constant on an entire usable endpoint cell \(U\), and the particular
pad mask certifying that term contains the indicated vertex bearing its tag.
If the group, label, and inverse-private-weight counts on a path are
bounded by \(\exp(C_qk^C)\), H's chosen mismatch accuracy gives
\begin{equation}
 \left\|\sum_Gh_G\right\|_2^2\le C\sum_G\|h_G\|_2^2.
 \label{eq:h-usable-synthesis}
\end{equation}
Labels and cells may be repeated at arbitrary endpoint times.  The
constant does not count these repetitions or the depth of a future
certificate.  If an output is restricted to a union \(E\) of finer
cells, then \((\mathbf1_E P_U)^*=P_U\mathbf1_E\); its adjoint range
is still the column span on the whole \(U\).  In particular, for a
common depth \(t\), certified unions \(E_G\) of \(t\)-cells, and
individual ball projections \(L_{G,t}\), duality gives
\begin{equation}
 \sum_G\|\mathbf1_{E_G}L_{G,t}h\|_2^2\le C\|h\|_2^2
 \quad\text{for one shared input }h.
 \label{eq:h-usable-rows}
\end{equation}

The binary localization proof in H now applies to the one current
tuple \((z_{0,s},\ldots,z_{3,s})\) on each original cell in
\(\mathcal A_R^{\rm main}\).  Each pad in a smeared filter acts on
that same operand.  The closure tests are \eqref{eq:h-closure-test},
the mismatch estimates concern fixed column tags, and the other
factors use only local \(L^2\) sizes.  The
\emph{Localization into current groups} Proposition of
\cite[Section~6]{H}, Equation~\Heq{14}, therefore gives at most
\(\operatorname{poly}(k)\) fixed search leaves and, for every
\(I\in\mathcal A_R^{\rm main}\),
\begin{equation}
 \begin{gathered}
 H_I(z_{\cdot,s})=
   \sum_\lambda\ \sum_{G\ {\rm current}}
     H_I(Z^\lambda_{0,G,s},\ldots,Z^\lambda_{3,G,s})+e_I,\\
 Z^\lambda_{l,G,s}=\overline L^\lambda_{l,G,s}
                         \mathcal B^\lambda_{l,s}z_{l,s},\qquad
 \sum_{I\in\mathcal A_R^{\rm main}}|I||e_I|\le Ce^{-k}|R|.
 \end{gathered}
 \label{eq:h-localization}
\end{equation}
Here \(\overline L\) is one positive average of ball projections,
and \(\mathcal B\) is a source-ordered history of cumulative filters
and complements of length \(O(\log\mathcal L)\), independent of
\(G\).  The history and slot choices are fixed by the search leaf
throughout the tree.  Only groups with an actual current central
use on \(I\) are summed in each leaf, and that sum may be empty.
The sum over groups stays inside
the contribution of the original cell; subadditivity distributes
the minimum with \(\eta\) only over the fixed leaves.

\subsection{\texorpdfstring{Variation of the localized \(B\) output}{Variation of the localized B output}}

Fix one slot, one search leaf, and its color, and suppress their
subscripts.  Write the full basic expression as
\[
 B_s=\sum_{\tau=1}^{h_0}\varepsilon_\tau Q_s^\tau F_s,
 \qquad h_0\le2,
\]
with its signs and constituent basic projection families fixed.
Define the full endpoint output
\[
 Z_{G,s}=\overline L_{G,s}\mathcal B_sB_s
\]
using the source positive ball average and the fixed
group-independent history of the leaf.  We evaluate this expression
on every needed endpoint of the complete inherited families within
the attempt, including endpoints where \(G\) has no current central
use.  All individual pad families and unused-cell continuations are
fixed for the calculation.  The same pad choice is coupled across
compared depths in each average.

We also use a shared output \(Z_s^{\rm sh}\) from H's
group-independent current-depth coarse histories on \(F_s\).
These are the histories of the fixed source leaf and the remaining
input-side histories obtained when a source transfer step removes
its prescribed positive output-side filters in decreasing rank.
They retain the source order, \(O(\log\mathcal L)\) main factors,
normalized positive pad averages or controlled complements, and at
most two fixed signed basic constituents.  Their complement expansions
have the source bound \(2^{O(\log\mathcal L)}\).  The identity is included.
The output has no outer ball and is counted once across its uses.

\begin{lemma}[Certified variation with depth inputs]
\label{lem:h-certified-variation}
For the full endpoint outputs
\(Z_{G,s}=\overline L_{G,s}\mathcal B_sB_s\) just defined, let
\(\mathcal J\) be a common finite collection of depth intervals
\([r,t]\) with edge-overlap multiplicity at most \(a\).
For each \(G,[r,t]\), let \(E_{G,r,t}\) be a union of \(t\)-cells
having a downstream actual-use certificate for \(G\).  Then
\begin{equation}
 \sum_{G,[r,t]\in\mathcal J}
   \|\mathbf1_{E_{G,r,t}}(Z_{G,t}-Z_{G,r})\|_2^2
       \le aG_0|R|.
 \label{eq:h-certified-variation}
\end{equation}
For one shared output \(Z_s^{\rm sh}\) in the stated source class,
let \(E_{r,t}\) be the union of the required \(t\)-cells across all
its uses.  Then
\[
 \sum_{[r,t]\in\mathcal J}
   \|\mathbf1_{E_{r,t}}(Z_t^{\rm sh}-Z_r^{\rm sh})\|_2^2
       \le aG_0|R|.
\]
This shared assertion has no group sum or group certificate.
All outputs before the attempt root are zero.
\end{lemma}

\begin{proof}
Expand the complements, choose one basic constituent, and choose
one individual pad in each average, coupled at all depths.  One
term of the full ball output is
\begin{equation}
 \widehat Z_{G,s}=L_{G,s}B'_sF_s,\qquad
 B'_s=P_s^m\cdots P_s^1.
 \label{eq:h-current-product}
\end{equation}
Here \(L_{G,s}\) is an individual nested ball projection and the
\(P^j_s\) are nested main projection families shared across groups.
Every nonzero \(B\) constituent has its specified basic projection
as the input-side factor \(P_s^1\).  An empty product is allowed for
a shared history.

Color the intervals into at most \(a\) collections with disjoint
interiors.  Within one color, omit zero intervals and order the
others as \(r_1<t_1\le r_2<t_2\le\cdots\).  Initially suppose both
endpoints lie within the attempt.  The exact difference of the
coupled term is
\begin{align}
 \widehat Z_{G,t}-\widehat Z_{G,r}
 ={}&(L_{G,t}-L_{G,r})B'_rF_r
       +L_{G,t}(B'_t-B'_r)F_r \notag\\
    &+L_{G,t}B'_t(F_t-F_r).
 \label{eq:h-variation-decomposition}
\end{align}
We first establish ordinary Bessel bounds for the two
operator-change terms, then evaluate their rows on \(F_r\).  The
last term changes only the shared input at the right endpoint.

For the first term, initially omit \(B'_r\).  On each \(r\)-cell
\(U\), split off
\[
 C_{G,r,t,U}=\mathbf1_{E_{G,r,t}\cap U}
                 (L_{G,t}-L_{G,r})\mathbf1_U.
\]
At each interval \([r,t]\), locality on the \(r\) partition makes the outputs of these split
rows disjoint.  For a fixed \(G\), the unrestricted increments on
the disjoint intervals are orthogonal, and output restriction is a
contraction, so the split rows have an ordinary Bessel bound.
Their adjoints take values in the sums of the whole ball ranges on
the \(r\)- and \(t\)-cells.  For a nonzero split row, a certified
\(t\)-cell in \(E_{G,r,t}\cap U\) supplies a downstream use below
\(U\).  It certifies \(U\) for its particular pad mask whenever
that mask is nonempty; an empty mask contributes the zero
projection.  Thus both endpoint ranges are eligible for
\eqref{eq:h-usable-synthesis}.  Applying that bound to the adjoint
sums makes these ordinary rows jointly Bessel across groups and
intervals.

Insert \(B'_r\) at the decision cell \(U\) by
\eqref{eq:h-compression}.  For a nonempty product the row adjoint
has the order
\[
 P_U^1\cdots P_U^m(L_{G,t}-L_{G,r})
                         \mathbf1_{E_{G,r,t}\cap U}.
\]
Its final action is \(P_U^1\), so its range is a main-catalog
combination on the entire \(U\).  The input-support and observation
depths are both \(r\), and the read index is \(r\).  With an empty
product, both endpoint ranges are combinations on the \(t\)-cells:
the earlier range has this description there by inheritance.
The input-support depth is then \(r\), the observation depth is
\(t\), and the read remains \(r\).  In both cases the next
input-support partition refines the current observation partition,
because \(t\le r_{\rm next}\).  Equation~\eqref{eq:h-row-use}, with
budget \(D_{\rm m}\), now bounds the squared sum of the first term
on the varying reads \(F_r\) by \(G_0|R|\).

For the second term, \eqref{eq:h-usable-rows} first bounds the
outer ball rows at each common right endpoint on an arbitrary
shared operand.  It remains to bound the ordinary rows from
\[
 B'_t-B'_r
 =\sum_{j=1}^m P_t^m\cdots P_t^{j+1}
              (P_t^j-P_r^j)P_r^{j-1}\cdots P_r^1.
\]
After the triangle inequality over \(j\), the left factors are
contractions.  For each \(j\), the middle increments on the disjoint intervals are
Bessel.  Split them into \(r\)-cell input-support rows; the exact
split-norm identity in the proof of Corollary~\ref{cor:row-projections}
preserves the Bessel bound.  Insert the right product by Compression.
If that product is nonempty,
the final adjoint range is a combination on its whole \(r\)-cell;
if it is empty, as it is for \(j=1\), inheritance represents both
increment ranges on the \(t\)-cells.  The respective observation
depths are \(r\) and \(t\), the input-support depth and read index
are \(r\), and the next input-support depth is at least \(t\).
We have obtained the ordinary Bessel bounds before applying
\eqref{eq:h-row-use} to the reads \(F_r\).  The resulting squared
sum is \(G_0|R|\).  This argument also bounds the operator change
for a shared no-ball history; an empty history has zero operator
change.

For the final term, the product is shared across groups, so the
endpoint Bessel bound gives
\[
 \sum_G\|\mathbf1_{E_{G,r,t}}L_{G,t}B'_t(F_t-F_r)\|_2^2
 \le C\|B'_t(F_t-F_r)\|_2^2
 \le C\|\mathbf1_R(F_t-F_r)\|_2^2.
\]
Equation~\eqref{eq:h-input-variation} sums these differences.
For a shared history the same proof bounds its full \(L^2(R)\)
difference before the single restriction to \(E_{r,t}\); for the
shared identity it is precisely the input-difference estimate.

At most \(a\) intervals cross the start of the attempt.  The zero
convention, \eqref{eq:h-usable-rows}, contraction, and
\(\|\mathbf1_R F_t\|_2^2\le|R|\) bound those terms.  Finally restore the
positive averages by Jensen's inequality with coupled pads and
restore the fixed signed constituents and complement expansions
by the triangle inequality in the finite Hilbert sum of rows.
Their source-controlled lengths are \(O(\log\mathcal L)\), so
their costs and the row losses fit \(G_0\).  Summing the interval
colors proves both assertions.
\end{proof}

\begin{corollary}[The endpoint histories used by filter transfer]
\label{cor:h-transfer-endpoint-variation}
Fix one source search leaf, one step in the
\emph{Outgoing filters and near-band defects} Lemma of
\cite[Section~6]{H}, and one choice in that step's expansion of
the complementary factors.  In its outgoing complementary slot
write \(x_{c,s}=\mathcal C_{c,s}B_{c,s}\) for exactly the
input-side source history left after deleting the outgoing positive
filter.  This inner history is shared across groups.  It retains
the source order, its normalized averages or
identities, its \(O(\log\mathcal L)\) length, and the specified
input-side basic \(Q_s^\tau\) in each nonzero constituent.  Write
\(x_{d,s}\) for the whole remaining source history in the other
complementary slot.  All endpoint evaluations use the complete
inherited families, with coupled pads at compared depths, and may
occur when the group is not current.

Let \(Q_{c,p,s}\), \(1\le p\le P\), be exactly that step's
individual outgoing projections, and put
\(Q_{c,0,s}=0\), \(Q_{c,P+1,s}=1\).
For a ball transfer, each \(Q_{c,p,s}x_{c,s}\) with
\(1\le p\le P\) is an eligible group-indexed ball endpoint, while
\(x_{c,s}\) is shared.  For a cumulative transfer the outer ball
has already been removed and all these endpoints are shared.
The other endpoint \(x_{d,s}\) either has its one source positive
ball average outermost, with shared inner history, or is shared
after that ball has been removed.

Each specified eligible ball output, using its source positive
average or one individual allowed pad as a singleton average,
satisfies \eqref{eq:h-certified-variation}.  Each specified shared
output satisfies the shared conclusion of
Lemma~\ref{lem:h-certified-variation}, with one union of required
right-endpoint cells across its uses.
\end{corollary}

\begin{proof}
The calculation for \(\widehat Z\) in
\eqref{eq:h-variation-decomposition} applies to each of these
remaining source histories before recombination: it uses their
original order and inherited families, not the removed output-side
factors.  For a no-ball history use the shared part of that proof.
The controlled signed constituents and averages are restored in
the same way.  This gives the asserted bound per history or
individual endpoint.

In particular, a nonterminal outgoing band
\((Q_{c,p,s}-Q_{c,p-1,s})x_{c,s}\) is a difference of two eligible
endpoints.  A terminal ball band is the shared \(x_{c,s}\) minus
the positive endpoint \(Q_{c,P,s}x_{c,s}\); the shared term is
counted once.  Splitting one band costs a fixed factor, and summing
the \(P+1\) bands retains the allowed polynomial cost in \(P\).
\end{proof}

\subsection{The input in a frozen mode row}

We next replace the fixed-input evaluation in
\emph{Composition with the lift and main histories} of
\cite[Section~8]{H}.  The modes in that construction form a
catalog separate from the main columns.  Let \(D_*\) be its budget,
and let \(K\) be the positive dyadic integer chosen by H for the
analysis-block length after its mode and gap parameters.  The
quantitative conventions are
\begin{equation}
 \begin{gathered}
 \log D_{\rm base},\ \log D_{\rm m}\le C_qk^C,\qquad
 \log D_*\le C_q(1+k+P)^C,\\
 K\le C_q(1+k+P)^{C_5},\qquad
 G_0=\exp\!\bigl(C_q(1+\log(2k))^C\bigr).
 \end{gathered}
 \label{eq:h-budget-hierarchy}
\end{equation}
The budget \(D_*\ge e^k\) covers mode counts, coefficient sums,
rational heights, and inverse private weights.  It also bounds the
available mode labels over relevant groups on a path and the actual
cell/group uses.  The horizontal parameter dimension of a tag or
of two tags together is at most
\(d_{\max}\le C_qk^C\).  All displayed polynomial degrees have
absolute upper bounds before \(q\) is chosen.

Partition the attempt depths into \([b,b+K)\), and put
\[
 u_b=b-4K,\qquad v_b=b-2K,\qquad t_b=b+K.
\]
Write \(r_s=2^{-s}|J|\) for the interval length at original depth \(s\).
In a \(B\) slot freeze the full endpoint output \(Z_{l,G,u_b}\),
with value zero when \(u_b\) precedes the attempt.  It is computed
from \(F_{l,u_b}\) on whole \(u_b\)-cells, whether or not the group
is current there.  A \(v_b\)-cell is retained for this operation
only when a descendant actual use occurs in the block.

\begin{figure}[ht]
\centering
\begin{tikzpicture}[x=0.95cm,y=1cm,every node/.style={font=\small}]
  \draw[-{Latex[length=2mm]}] (0,0.8)--(11.7,0.8) node[right]{depth};
  \foreach \x in {0.5,4.5,8.5,10.5} {\draw (\x,0.66)--(\x,0.94);}
  \draw[fill=black!8] (8.5,1.1) rectangle (10.5,1.6);
  \draw[densely dashed] (9.5,0.8)--(9.5,1.6);
  \node[fill=black!8,inner sep=1.5pt] at (9.5,1.35) {$s$};
  \node at (9.5,1.86) {actual uses};
  \draw[{Latex[length=1.7mm]}-{Latex[length=1.7mm]}] (0.65,2.25)--(8.35,2.25);
  \node[fill=white,inner sep=2pt] at (4.5,2.25) {$4K$ depths};

  \node[anchor=north,align=center] at (0.5,0.45)
       {$u_b=b-4K$\\freeze $Z_{l,G,u_b}$};
  \node[anchor=north,align=center] at (4.5,0.45)
       {$v_b=b-2K$\\common test cell};
  \node[anchor=north,align=center] at (8.5,0.45)
       {$b$\\block begins};
  \node[anchor=north,align=center] at (10.5,0.45)
       {$t_b=b+K$\\block ends};
\end{tikzpicture}

\caption{The depth convention of the first-match construction in
\cite[Section~8]{H}.  In the analysis block \([b,b+K)\), the basic
output in slot \(l\) is frozen at \(u_b=b-4K\), tested on common
cells at \(v_b=b-2K\), and used at actual depths
\(b\le s<t_b=b+K\).  The deterministic input block partitions in
Theorem~\ref{thm:depth-local} need not align with these analysis blocks.}
\label{fig:frozen-depths}
\end{figure}

A primary in this mode construction is an individual main column
\(\Phi\) matched to the maximal pad of the role's outer ball
filter.  At its first match on a cell \(I_0\), the
\emph{First-match expansion and bounded lift} Lemma of
\cite[Section~8]{H} chooses a finite list of augmented modes.
For each prospective branch offset \(\rho\), the modes use H's
cutoff \(\zeta^\rho\) and gauge \(p^\rho\) in the form
\[
 \psi_m(y)=
 \left(\zeta^\rho(y)e(\kappa_l p^\rho(y))
       e(\lambda_m y^2+\gamma_m y),
       \sqrt{\epsilon_{\rm m}}\,\mathbf e_m\right).
\]
The list, its labels, and its scalar expansion coefficients
\(a_{\Phi,m}\) are retained on every relevant descendant of the
first match; separate first-match branches may use separate private
labels.  They are not resampled at later freezing times, including
when a coefficient subsequently frozen is zero.  All modes of
one primary have one \(\lambda_m\).  The quantities
\(|\lambda_m||I_0|^2\), \(\sum_m|a_{\Phi,m}|\), their number, and
the heights of their relative linear-frequency vectors are bounded
by a quantity whose logarithm is polynomial in \(P,k\).

For a fixed prospective offset the same source Lemma gives one
pointwise bounded linear map \(\mathcal R_{l,G}\), fixed across
blocks.  On an outer main output
\(\sum_\Phi\beta_\Phi\Phi\) with coefficients constant on an
endpoint cell, it gives exactly
\(\sum_{\Phi,m}\beta_\Phi a_{\Phi,m}\psi_m\).  Thus
\[
 \widetilde Z_{l,G,b}=\mathcal R_{l,G}Z_{l,G,u_b}
\]
has coefficients constant on its whole \(u_b\)-cell.  The mode
labels and their private coordinates use \(D_*\); the main
endpoint projections still use \(D_{\rm m}\).

Write \(\mathcal H_{\rm mode}\) for the augmented mode Hilbert space.

\begin{lemma}[Frozen composition with depth inputs]
\label{lem:h-frozen-composition}
Fix a slot, one search leaf and its color, and a prospective choice
of the offsets for its groups, with the corresponding block-independent
maps \(\mathcal R_{l,G}\).  For every group \(G\) of that color, let \(P_e\) be finitely many
local mode projections: \(P_e\) is the orthogonal projection in
\(L^2(J_e;\mathcal H_{\rm mode})\) onto constant-coefficient combinations
of the restrictions of its available augmented modes.  Extend it to
\(L^2(R;\mathcal H_{\rm mode})\) by restricting the input to \(J_e\)
and extending the output by zero.  Suppose these rows have joint squared
Bessel bound at most \(G_0\) in each group.

Suppose \(J_e\subseteq U_e\), where \(U_e\) is a main endpoint
cell with a downstream actual-use certificate and is usable for
each nonzero pad range contributing to
\(\overline L_{l,G,U_e}\).  The row's mode labels are available
already at \(U_e\).  Write
\[
 B_{l,U_e}=\sum_{\tau=1}^{h_0}\varepsilon_\tau
       Q^\tau_{l,U_e}F_{l,s(U_e)},\qquad h_0\le2,
\]
for its specified basic-\(B\) expression, and let
\(\mathcal B_{l,U_e}\) be the shared main history of the leaf.
Then
\begin{equation}
 \sum_{G,e}
 \|P_e\mathcal R_{l,G}\overline L_{l,G,U_e}
                  \mathcal B_{l,U_e}B_{l,U_e}\|_2^2
       \le G_0|R|.
 \label{eq:h-frozen-composition}
\end{equation}
Endpoint depths may be arbitrary and need not be current; any
number of rows may share one endpoint cell.  The bound is uniform
in the prospective choices of offsets.
\end{lemma}

\begin{proof}
We first form an ordinary joint Bessel operator using the fixed
lift, outer main projections, and usable-endpoint synthesis.
We then insert the shared histories and the constituent basic
projection, and only afterward evaluate the rows on their varying
endpoint inputs.

For a fixed group the rows \(P_e\mathcal R_{l,G}\) are Bessel
because \(\mathcal R_{l,G}\) is one fixed bounded operator.  They
are local to \(J_e\), hence to the containing decision cell
\(U_e\).  Compression inserts an outer individual main projection
on \(U_e\); the triangle inequality in the finite Hilbert sum of
rows restores its positive average.  If \(A_G\) denotes this row
operator, an individual row adjoint has the order
\[
 L_{l,G,U_e}\mathcal R_{l,G}^*P_e.
\]
Its range is the main ball range on the whole \(U_e\).  The pad
usability hypothesis and \eqref{eq:h-usable-synthesis} therefore
give
\[
 \left\|\sum_G A_G^*w_G\right\|_2^2
 \le C\sum_G\|A_G^*w_G\|_2^2
 \le G_0\sum_G\|w_G\|^2.
\]
By duality, the ordinary rows are jointly Bessel across all groups
and rows.  Their use of the whole endpoint \(U_e\) is the reason
the synthesis bound applies despite the smaller output cell \(J_e\).

Apply Compression again at the decision cells \(U_e\) to insert
the shared histories and one fixed basic constituent.  Expand
complements into their signed ordered products and retain
normalized positive averages.  For each coupled expanded term,
the resulting ordinary rows \(T_e\) satisfy
\begin{equation}
 \sum_{G,e}\|T_e h\|_2^2\le G_0\|h\|_2^2
 \quad\text{for every }h\in L^2(R;\mathcal H).
 \label{eq:h-frozen-operator}
\end{equation}
Every nonzero basic constituent has its input-side projection
\(Q^\tau_{l,U_e}\), even if all other selected factors are
identities.  Thus the final action of \(T_e^*\) is
\(Q^\tau_{l,U_e}\), and its range is a main-catalog combination
on the entire \(U_e\).  The composed row has input support on
\(U_e\).  Grouping rows by \(s(U_e)\), the full depth partition
serves as both input-support and observation partition, and the
read is \(F_{l,s(U_e)}\).  These partitions refine and the read
indices increase.  Equation~\eqref{eq:h-row-use}, applied to
\eqref{eq:h-frozen-operator} with budget \(D_{\rm m}\), proves
\eqref{eq:h-frozen-composition} after restoring the controlled
signed terms and averages.

The mode projections entered the ordinary Bessel operator before
this last evaluation.  The final adjoint range is the basic main
range on \(U_e\), so the row estimate uses \(D_{\rm m}\).  For a
row from block \(b\), its read and observation depth are
\(s(U_e)=u_b\), regardless of the later mode-output cell \(J_e\)
or the actual use depth.
\end{proof}

Lemmas~\ref{lem:h-certified-variation} and
\ref{lem:h-frozen-composition} now provide the cross-depth
variation and frozen-evaluation bounds needed for H's mode split.
We next verify the resulting side energy and component data.

\subsection{Side and core data from the frozen modes}

Fix a prospective choice of the group offsets.  On each retained \(v_b\)-cell, make
H's permitted choice to keep the full available mode list frozen
at \(u_b\), including every matched primary and every
zero-coefficient label.  PRE and CURRENT mean exactly the graphs
of Equation~\Heq{22} in \cite[Section~8]{H} on that list and that
same \(v_b\)-cell.  Their edge tests use the supremum of the
affine-slope difference
\(\omega_m(x)-\omega_{m'}(x)\), where
\(\omega_m(x)=2\lambda_mx+\gamma_m\), after the source's permitted
integral shifts.  The comparison scales are \(r_{v_b}\) for PRE
and \(r_{t_b}\) for CURRENT.  A minimal PRE component and a minimal
CURRENT component are their respective height-zero components.
Available lists and old labels persist on relevant descendants;
admitted edges persist or grow as the source height or descendant
block advances.

We retain H's source-defined good predicates and core/side
classification.  The \emph{Coherence of good predicates} Lemma
says that a minimal PRE component classified as side consists
entirely of side modes.  It also says that the upper good-height
masks select unions of complete core subtotals inside minimal
CURRENT components.  A minimal CURRENT component may contain
side modes as well.  At an upper height, a successful CURRENT
component means one containing a core mode good at that height;
its good modes have one PRE predecessor at the prescribed lower
height.

For a \(B\) role, let \(V_{l,G,b}\) be the exact subtotal of the
side modes of \(\widetilde Z_{l,G,b}\), including their private
coordinates, on retained \(v_b\)-cells.  Set it to zero on branches
without a needed output or certificate.  For the fixed prospective
offset it is one fixed function throughout the block.  We verify
the source square-energy estimate for these new frozen operands.

An event in the proof of \emph{Square energy of the linear side}
in \cite[Section~8]{H} consists of a block, a retained \(v_b\)-cell
\(J\), and a minimal PRE component \(\mathcal E\) of side modes.
The \emph{Separated mode rows} Lemma and the
\emph{Bad-density chains} Lemma organize its projection rows into
ten block congruence classes and
\(O(d_{\max}\log(2D_{\rm m}))\) chain levels.  The source chain
calculation absorbs its larger mode-height parameter before this
final length bound.  Thus the per-group squared Bessel cost is
\(G_0\).  Lemma~\ref{lem:h-frozen-composition}, with \(U_e\) the
containing \(u_b\)-cell, gives the source Equation~\Heq{24}:
\[
 \sum_{G,b,J,\mathcal E}
       \|P_{\mathcal E,J}\widetilde Z_{l,G,b}\|_2^2
       \le G_0|R|.
\]
Here \(P_{\mathcal E,J}\) is the local projection onto that
component's augmented modes.  If \(T_{\mathcal E,J}\) is its
exact coefficient subtotal, the source separation and the bound
\(\sum_m|\alpha_m|\le D_*\) on the frozen cell give
\[
 \|P_{\mathcal E,J}\widetilde Z-T_{\mathcal E,J}\|_{2,J}
       \le D_*^{-190}.
\]
The path budgets sum these additive errors.  Almost orthogonality
of the distinct minimal PRE subtotals then yields
\begin{equation}
 \sum_{G,b}\|V_{l,G,b}\|_2^2\le G_0|R|.
 \label{eq:h-linear-energy}
\end{equation}
This deduction is uniform in the prospective choices of offsets.  It uses the
persistent lists and graphs and each endpoint's own coefficient
subtotal; it requires no relation between coefficients frozen in
different blocks.

Let \(A_{l,G,b}\) be the exact core subtotal.  In a \(B\) role
define
\[
 S^{\rm roll}_{l,G,s}=Z_{l,G,s}-Z_{l,G,u_b},\qquad
 S_{l,G,s}=S^{\rm roll}_{l,G,s}+V_{l,G,b},
 \quad b\le s<b+K.
\]
The physical identity in
\emph{Fixing offsets and separating the side terms} of
\cite[Section~8]{H} is
\begin{equation}
 \begin{gathered}
 (Z_{l,G,s})_{\rm phys}
   =(A_{l,G,b})_{\rm phys}+(S_{l,G,s})_{\rm phys}
          +\mathcal E_{l,G,b},\\
 \mathcal E_{l,G,b}
   =(Z_{l,G,u_b})_{\rm phys}
          -(\widetilde Z_{l,G,b})_{\rm phys}.
 \end{gathered}
 \label{eq:h-core-side}
\end{equation}
In a basic-\(U\) role there is no core: set
\(S_{l,G,s}=Z_{l,G,s}^{\,U}\) and take no physical error.
The core and side subtotals retain their augmented coordinates
in all norm estimates.

This decomposition determines the remaining estimates.  Terms
with at least three side positions will use their path-square
bounds at a power \(p\in(2/3,1)\).  Terms with exactly two sides
will restore the two whole current complements and use filter
transfer.  Terms with at least three cores will use H's
conditional counting estimate.  For a powered within-cell group
contribution \(x\), the truncation gives
\(\min(\eta,x)\le\eta^{1-p}x^p\).  The transfer's bands with
adjacent pad indices will be estimated at power one, with their
prescribed fixed-offset sums kept inside the cell.

It remains to select the offsets.  The main catalogs, the
coefficients \(\beta_\Phi\) in
\(Z_{l,G,u_b}=\sum_\Phi\beta_\Phi\Phi\) on each whole \(u_b\)-cell,
and all actual main uses were fixed before the prospective
offsets.  With the source accuracy \(\delta_c=\exp(-k^{C_c})\),
where \(C_c\) is chosen after the main budgets, H's explicit lift
and offset average give
\[
 \mathbb E_\rho\sum_{l,G,I\ {\rm used}}|I|
       \|\mathcal E_{l,G,b(I)}\|_{2,I}^2
 \le D_{\rm m}^C d_{\max}\delta_c\,|R|.
\]
This sum is over all those pre-existing actual uses, before any
offset-dependent cuts.  The coefficient and physical supremum
bounds are on the whole \(u_b\)-cell.  By
\eqref{eq:h-input-size}, its operand \(F_{l,u_b}\) has normalized
norm at most one, so the lag \(4K\) introduces no size factor.
Choose offsets satisfying the displayed bound, with the
inverse-main-budget accuracy \(\delta_c\) prescribed by H.
The uniform estimate \eqref{eq:h-linear-energy} remains valid
for this choice, and subsequent pruning only removes terms from
the nonnegative error sum.

The chosen whole augmented states \(S_{l,G,s}\) are now fixed.
Realize H's helper closure on these states and their named
histories by the same fixed-base calculation
\eqref{eq:h-closure-test}--\eqref{eq:h-closure-prefix}.  There
are at most \(\exp(C_qk^C)\) such bases per active cell over the
specified slots, groups, and patterns.  They are whole states,
with no choice of a mode, height, or amplitude as a helper base.
Their existing private coordinates are distinct from the fresh
helper directions.  This realizes the anticipated helper count
within \(D_{\rm m}\).  The dependent size and counting cuts
described next are imposed after these choices.

The separate integrated square bounds for the side components
follow from the two new lemmas and the basic energy.  For the
rolling component, apply \eqref{eq:h-certified-variation} to
\([u_b,s]\), \(b\le s<b+K\), with the unions of actual \(s\)-cells
as endpoint sets; these intervals have \(O(K)\) overlap.  At each
depth the use cells for fixed \(G,b\) are disjoint, so a fixed
\(V_{l,G,b}\) is counted at most \(K\) times.  In a \(U\) role,
\eqref{eq:h-usable-rows} at each depth and the basic energy in
\eqref{eq:h-base-data} bound the current outputs.  Consequently
\begin{equation}
 \begin{aligned}
 \sum_{s,G}\sum_{I\ {\rm used\ at}\ s}|I|
       \|S^{\rm roll}_{l,G,s}\|_{2,I}^2
       &\le CK G_0|R| &&\text{in a \(B\) role},\\
 \sum_{s,G}\sum_{I\ {\rm used\ at}\ s}|I|
       \|V_{l,G,b(s)}\|_{2,I}^2
       &\le K G_0|R| &&\text{in a \(B\) role},\\
 \sum_{s,G}\sum_{I\ {\rm used\ at}\ s}|I|
       \|S_{l,G,s}\|_{2,I}^2
       &\le P^C|R| &&\text{in either role},\\
 \sum_{s,G}\sum_{I\ {\rm used\ at}\ s}|I|
       \|Z^{\,U}_{l,G,s}\|_{2,I}^2
       &\le G_0|R| &&\text{in a \(U\) role}.
 \end{aligned}
 \label{eq:h-side-square}
\end{equation}
The first two inequalities concern a \(B\) role.  Their sum
bounds its whole \(S\) by the squared triangle inequality.
The fixed polynomial in \(k\) from the basic energy is included
in \(G_0\).

H's componentwise first-crossing construction in
\cite[Section~7, Side estimate]{H} applies to this finite list of
components.  Apply it separately to the nonnegative partial path
sums for \(S^{\rm roll}\), for the repeated \(V\), and for the
basic-\(U\) output when present, using their respective integrated
bounds above.  By increasing the thresholds by fixed factors,
their first-crossing cells have any prescribed small fixed total
relative length.  On the retained paths this gives
\begin{equation}
 \begin{aligned}
 \sum_{s,G}\|S^{\rm roll}_{l,G,s}\|_{2,I_s(x)}^2
       &\le P^C &&\text{in a \(B\) role},\\
 \sum_{s,G}\|V_{l,G,b(s)}\|_{2,I_s(x)}^2
       &\le P^C &&\text{in a \(B\) role},\\
 \sum_{s,G}\|Z^{\,U}_{l,G,s}\|_{2,I_s(x)}^2
       &\le P^C &&\text{in a \(U\) role},\\
 \sum_{s,G}\|S_{l,G,s}\|_{2,I_s(x)}^2
       &\le P^C &&\text{in either role}.
 \end{aligned}
 \label{eq:h-side-cap}
\end{equation}
Only actual uses are summed.  The final cap follows from the
separate component caps and
\(\|S^{\rm roll}+V\|^2\le2\|S^{\rm roll}\|^2+2\|V\|^2\).
The exponent may be increased for the polynomially many
configurations where a power cost is allowed.

These separate caps also give the local input for H's component
sizes.  At a retained \(B\) use, the current local size of \(Z_s\)
is bounded by contraction and \eqref{eq:h-base-data}.  The
bounded lift and \(Z_{u_b}=Z_s-S^{\rm roll}_s\) therefore give
\[
 \|\widetilde Z_{l,G,b}\|_{2,I}+\|V_{l,G,b}\|_{2,I}
 \le C\bigl(\|Z_{l,G,s}\|_{2,I}
             +\|S^{\rm roll}_{l,G,s}\|_{2,I}\bigr)
             +\|V_{l,G,b}\|_{2,I}
 \le P^C.
\]
The core subtotal \(A=\widetilde Z-V\) has the same type of local
bound.  H's local \(L^2\) form estimate and Cauchy--Schwarz over
the original actual uses now turn the selected integrated
physical-error bound into total absolute error \(O(e^{-k}|R|)\).

We next record the representative and top-count inputs for the
core calculation.  A successful component's type is the pair of
floor logarithms, to base \(1.1\), of its unique good PRE
predecessor's primary count and mode count.  For a fixed role,
group, upper height, and type, the
\emph{Hereditary representatives} Lemma of
\cite[Section~8]{H} chooses a mode label \(a\) in each such
predecessor and inherits it on later contained predecessors.
Write \(z\) for the group's primary center and \(\theta_z\) for
the horizontal vector in its tag.  H uses
\(\vartheta=\theta_z/Q_0\), where \(Q_0\) is H's denominator-clearing
integer common to the roles of this group.  Write \(L_h,E_h,\Xi\) for the
source shift denominator, shift range, and separation threshold.
For distinct representatives on comparable actual cells \(I,I'\)
and \(x_*\in I\cap I'\), it gives
\begin{equation}
 |\omega_a(x_*)-\omega_{a'}(x_*)+n\cdot\vartheta/L_h|
   >\frac{\Xi}{2}\max(|I|^{-1},|I'|^{-1})
 \quad(n\in\mathbb Z^{\dim\theta_z},\ |n|_\infty\le E_h).
 \label{eq:h-representative-separation}
\end{equation}
The identification across depths is the persistent mode label
and its value \(\omega_a(x_*)\).

For an actual cell \(I\) and a successful component
\(\mathcal F\), H uses
\[
 \sigma_{\mathcal F}^2
 =\sum_{\substack{\mathcal D\subseteq\mathcal F\\
                    \mathcal D\ {\rm minimal\ CURRENT}}}
    \left(\|A_{\mathcal D}\|_{2,I}
       +D_*^{-30}\!\!\sum_{\substack{m\ {\rm core}\\m\in\mathcal D}}
                         |\alpha_m|\right)^2.
\]
Here \(A_{\mathcal D}\) is the complete core subtotal within
\(\mathcal D\); side labels in that component are kept separately.
The \emph{Component-size comparison} Lemma of
\cite[Section~8]{H} gives the same-cell bounds
\[
 \sigma_{\mathcal F}^2
 \le C\bigl(\|P_{\mathcal F,I}\widetilde Z_{l,G,b}\|_{2,I}^2
             +\|V_{\mathcal F}\|_{2,I}^2+D_*^{-50}\bigr),
 \qquad
 \sum_{\mathcal F}\|V_{\mathcal F}\|_{2,I}^2
       \le C\|V_{l,G,b}\|_{2,I}^2.
\]
The projection \(P_{\mathcal F,I}\) uses all mode labels of
\(\mathcal F\) on \(I\), and \(V_{\mathcal F}\) is their exact
side subtotal.  This is an almost-orthogonality estimate for
arbitrary coefficient subtotals on one cell.

For a fixed source role, upper height, type, and dyadic threshold
\(D_*^{-1}\le\sigma\le D_*^2\), H takes, for each group and
representative, the first actual-use cells on each spatial branch
for which \(\sigma_{\mathcal F}\ge\sigma\).  Several incomparable
tops for the same group and representative are allowed.  The
resulting top rows
\(P_{\mathcal F,I}\) are separated mode rows.  They are supported
on \(I\subseteq U_e\), the \(u_b\)-cell of the frozen list, and
the top itself is a downstream certificate.  Applying
Lemma~\ref{lem:h-frozen-composition} gives squared energy
\(G_0|R|\) for their frozen projections.  Summing the side
term of the size comparison first at one depth and then through
a block gives at most
\(K\sum_{G,b}\|V_{l,G,b}\|_2^2\le KG_0|R|\).  The
\(D_*^{-50}\) floor is absorbed because
\(\sigma^2\ge D_*^{-2}\).  Hence
\[
 \sigma^2\sum_{\rm tops}|I|\le C(1+K)G_0|R|.
\]
There are polynomially many choices in \(P,k\) of height, type,
and threshold.  The first-crossing construction of the
\emph{Top count} Proposition in \cite[Section~8]{H} gives, on
retained paths,
\begin{equation}
 \sum_G|\mathcal A_j^G|
       \le\min(D_*,P^C\sigma_j^{-2}).
 \label{eq:h-top-count}
\end{equation}
Here \(\mathcal A_j^G\) is the set of distinct representatives
actually used in the bin \([\sigma_j,2\sigma_j)\) along the retained
path in question, for the fixed role, height, and type.  The full
label budget supplies the
minimum for \(\sigma_j<D_*^{-1}\); bins above the tested range
are empty after the source budget enlargement.  Finally, the
local bounds on \(\widetilde Z\) and \(V\) above and the
component-size comparison give \(\sigma_{\mathcal F}\le P^C\)
on retained uses.  These are the local sizes and role counts
needed by the core calculation.

\subsection{The side estimate}

We use the \emph{Side estimate} Proposition of \cite[Section~7]{H}, with the
component bounds just verified.  Suppose first that three
positions have sizes whose path squares obey
\eqref{eq:h-side-cap}.  Along a path put \(a_s=(\sum_Ga_{s,G}^2)^{1/2}\)
and define \(b_s,c_s\) analogously, using a shared size directly
when appropriate.  Cauchy--Schwarz in the group index
bounds the absolute group sum by \(P^C a_sb_sc_s\), using the
retained bound \(P^C\) for the fourth local size.  For
\(2/3<p<1\), Hölder and \(3p>2\) give
\[
 \sum_s(a_sb_sc_s)^p
 \le
 \left(\sum_sa_s^2\right)^{p/2}
 \left(\sum_sb_s^2\right)^{p/2}
 \left(\sum_sc_s^2\right)^{p/2}.
\]
Integrating the path caps therefore gives total length-weighted
\(p\)th-power cost \(P^C|R|\).  The power is taken after the
group sum at each original cell.

For exactly two sides, replace their two core complements in
\eqref{eq:h-core-side} by the whole current \(Z\) arguments.
The corrections have a third side, apart from the physical errors
already estimated.  We recall the source filter algebra for these
whole complements.  Let the side positions be \(a,b\) and the
complementary positions be \(c,d\).  Expand each complementary
factor \(1-\overline Q\) as identity minus a positive average.
The \emph{Outgoing filters and near-band defects} Lemma of
\cite[Section~6]{H} removes the positive filters from \(c,d\)
in decreasing rank and transfers helper averages to the fixed
whole state \(S_a\).  Its final main term is
\[
 H_I(\mathcal D_{a,s}S_a,S_b,B_c,B_d),
\]
where \(\mathcal D_{a,s}\) has \(O(\log\mathcal L)\) helper
factors.

For one step from \(c\) to \(a\), write \(Q_{l,p}\),
\(1\le p\le P\), for the outgoing and corresponding helper
projections, and put
\[
 Q_{l,0}=0,\quad Q_{l,P+1}=1,\quad
 A_l^p=Q_{l,p}-Q_{l,p-1},\quad
 \overline Q_l=P^{-1}\sum_{p=1}^P Q_{l,p}.
\]
Here \(x_c\) is the exact remaining input-side source history
on \(B_c\), \(x_d\) is the whole remaining history on \(B_d\),
\(x_a\) is the already transferred helper history on \(S_a\),
and \(x_b=S_b\).  The exact difference is
\begin{equation}
 H_I(x_a,x_b,\overline Q_cx_c,x_d)
       -H_I(\overline Q_ax_a,x_b,x_c,x_d)
 =\sum_{p,q'=1}^{P+1}\frac{p-q'}P
       H_I(A_a^p x_a,x_b,A_c^{q'}x_c,x_d).
 \label{eq:h-transfer}
\end{equation}
The pad projections are nested on this cell, so the \(A_l^p\)
are orthogonal bands summing to the identity.  The average has
weight \((P-p+1)/P\) on band \(p\), which proves the identity.
H's closure and tag-separation estimates bound the sum with
\(|p-q'|>2\) by
\(D_{\rm base}^{-C_3}\|S_a\|_{2,I}\|S_b\|_{2,I}\), for any
prescribed fixed \(C_3\) after choosing its numerical accuracies.
The helper closure here uses the already fixed whole \(S_a\)
as its base in \eqref{eq:h-closure-test}.  For the retained terms,
keep the entire band sum at each fixed offset \(q'-p\).

In the final main term, write \(B_c=Y_c-U_c\) and
\(B_d=Y_d-U_d\).  Each correction has a third basic side.
The remaining complementary pair \(Y_c,Y_d\) has form norm at
most \(\eta\) against the two side inputs, by the full-pair
statement of Proposition~\ref{prop:hbase-data}.  Contraction of
the helper factors and Cauchy--Schwarz in the side sizes give
absolute main cost \(\eta P^C|R|\).  The numerical transfer
errors are summable with the chosen inverse-base-budget accuracy.

For a retained offset, write
\(C_{G,s}^{q'}=A_c^{q'}x_{c,s}\) for the outgoing complementary
band stack and \(D_{G,s}=x_{d,s}\) for the other complement.
Freeze both at the block start \(b\), using their full endpoint
evaluations there.  These endpoints differ from the mode
freezing depth \(u_b\).  Corollary~\ref{cor:h-transfer-endpoint-variation}
applies to exactly these source histories.  For a ball transfer,
a nonterminal band is the difference of two eligible positive
ball endpoints; its terminal band is the shared inner history
minus the positive \(P\)-pad endpoint.  For a cumulative transfer,
the outer ball has already been removed and all the complementary
endpoints are shared.  The other complement has its one source
positive outer ball or is shared.  Thus each band costs a fixed
multiple of its endpoint estimates, and the shared histories
are counted once across their uses.  Summing the \(P+1\) bands
and the \(O(K)\) overlap of \([b,s]\) retains the allowed
polynomial cost \(P^C|R|\).
Their nonnegative partial path sums enter the power-cost first-crossing
cuts.

At a required endpoint the pad bands are orthogonal, and the
remaining source histories are contractions of bounded basic
inputs.  Choose \(C_{\rm fr}\ge1\) to dominate the retained
current local norm of the band stack and of the other complement.
A freeze error has a third path-square size by the preceding
endpoint bounds.  If the other frozen complement is large,
write it as its bounded current value plus its difference.
Likewise, if a frozen local size exceeds \(2C_{\rm fr}\) while
the current one is at most \(C_{\rm fr}\), then the difference
exceeds \(C_{\rm fr}\), and the flag is at most
\(C_{\rm fr}^{-2}\) times the squared difference size.  Its square root is therefore
another path-square size.  Certified variation applies on the whole ancestor cells with downstream
certificates; propagation through a block costs at most \(K\)
in this power budget.  H's freeze, large-complement, and
restoration errors therefore have total \(p\)th-power cost
\(P^C|R|\).  A large-complement flag is charged to that error
sum; only a first crossing of its accumulated budget makes a
restart cut.  On the remaining uses the frozen complements have
bounded local sizes on their entire ancestor hulls in the block.

\subsubsection{Fixed stacks for the two side inputs}

Fix a source transfer step, a near offset, and a search leaf.
Under these fixed data, let \(\chi\) record the source auxiliary
side-component and binary-length names for the two side stacks.
Pairing their choices raises the auxiliary count to a fixed power,
so a common auxiliary list has at most \(G_0\) entries after enlarging \(G_0\) within
its stated class; unused combinations are set to zero.  Write
\(\nu=(G,b,\chi)\) for the full stack index.  Concrete binary
interval positions and band labels \(p\) are stack coordinates,
and individual pads remain in coupled averages.  For one full
index, a stored side stack is a function into a finite Hilbert
direct sum:
\[
 \begin{gathered}
 \mathrm{stack}_{\mathrm{untouched},\nu}(y)
       =(W_{\nu,i}(y))_{i\in\mathcal I_\nu}
       \in\bigoplus_{i\in\mathcal I_\nu}\mathcal H,\\
 \|\mathrm{stack}_{\mathrm{untouched},\nu}\|_2^2
       =\sum_i\|W_{\nu,i}\|_2^2.
 \end{gathered}
\]
The banded stack \(\mathrm{stack}_{a,\nu}\) has in addition the
band index \(p\in\{1,\ldots,P+1\}\).  Here \(\mathcal H\) is the
augmented ambient Hilbert space of the stored coordinates.  Each
whole coordinate is restricted, when formed, to the union of the
use cells selecting its position and valid band.  Later cuts
remove summands without changing that stored function.  At a use,
each side selects one position; the banded side retains all valid
matched band indices at the fixed offset.

The source stack construction gives the squared energies
\begin{equation}
 \sum_\nu\|\mathrm{stack}_{a,\nu}\|_2^2\le PG_0|R|,
 \qquad
 \sum_\nu\|\mathrm{stack}_{\mathrm{untouched},\nu}\|_2^2\le G_0|R|.
 \label{eq:h-stacks}
\end{equation}
Both sums range over all full indices \(\nu=(G,b,\chi)\).
The \(G_0\) cardinality overhead counts only the auxiliary labels
\(\chi\); both energies sum all groups and blocks.
Their quantitative purpose is the following.  The whole-hull
maximal and sparse calculation below turns these energies into
a root-size sum with cost \(P^{1/2}G_0|R|\).  The coefficient
\(O(P^{-1})\) in the matched transfer bands then leaves the
gain \(P^{-1/2}\).  We first check the fixed-vector arguments
behind \eqref{eq:h-stacks}.

For a sequence \(W_s\) in one block, write
\[
 W_s=W_b+\sum_{[r,t]\in\mathcal P(b,s)}(W_t-W_r),
\]
where \(\mathcal P(b,s)\) is the binary decomposition of the
prefix.  At a fixed length its intervals are disjoint and at
most one is selected at a use.  Store each resulting function
on its selecting-use union as above.  For a basic-\(U\) side,
store each already formed depth function directly;
\eqref{eq:h-side-square} bounds its energy, and subsequent
helper products are local contractions.

For a linear side \(V_{G,b}\), let \(T_s\) be the product of
the subsequent helper projections, including one endpoint of
a band difference when needed.  The products-of-prefixes
estimate of \cite[Section~3]{H}, on at most \(K+1\) times, gives
for each fixed binary length
\[
 \|T_bV_{G,b}\|_2^2+
 \sum_{[r,t]\ {\rm at\ that\ length}}
       \|(T_t-T_r)V_{G,b}\|_2^2
 \le (C\log(K+2))^{O(m)}\|V_{G,b}\|_2^2.
\]
The input \(V_{G,b}\) is the one fixed function on the block.
Use \eqref{eq:h-linear-energy}, sum the logarithmic length
choices, and pay at most \(P+1\) band indices for the first
stack.

For a rolling side put \(W_s=T_s(Z_{G,s}-Z_{G,u_b})\).
Its base value uses \eqref{eq:h-certified-variation} on
\([u_b,b]\), whose overlap across blocks is bounded.  Its
increments split exactly as
\[
 W_t-W_r=T_t(Z_{G,t}-Z_{G,r})
             +(T_t-T_r)(Z_{G,r}-Z_{G,u_b}).
\]
The first term uses certified variation on the whole endpoint
cells containing its selected use cells.  For the second, also
decompose \([u_b,r]\) into binary intervals.  Fix an outer
length, an inner length, and an inner interval \([a,v]\).
Every outer position selecting this interval acts on the one
fixed function
\[
 X=\mathbf1_\Omega(Z_{G,v}-Z_{G,a}),
\]
where \(\Omega\) is the union of its required \(r\)-cells.
Since \(v\le r\), locality of \(T_t-T_r\) gives the required
outputs from \(X\).  Enlarge \(\Omega\) to the containing whole
\(v\)-cells to bound its norm; those cells have downstream
certificates.  The fixed-vector product-jump estimate sums the
outer positions by
\((C\log(K+2))^{O(m)}\|X\|_2^2\).  Then
\eqref{eq:h-certified-variation} sums these squared inputs over
groups and the disjoint inner intervals at that length.  The
\(O(K)\)-depth windows for different blocks have bounded
overlap.  The two endpoint outputs in \(X\) use \(F_v\) and
\(F_a\), but \(X\) itself is fixed once \(a,v\) are fixed, as
required for the outer product-jump estimate.

This proves \eqref{eq:h-stacks}.  Coupled pad averages are restored
by Jensen's inequality.  When the fixed search-leaf and transfer-step
choices are restored later, their auxiliary overhead fits an enlarged
\(G_0\).

We collect the main-row and stack costs here.  For
\(m=O(\log\mathcal L)=O(\log k)\), the Compression factor in
\eqref{eq:h-compression} satisfies
\[
 \log\mathfrak C(m,D_{\rm m})
 \le C(m+1)\bigl(\log(m+1)+\log\log(2D_{\rm m})\bigr)
 \le C_q(1+\log(2k))^C.
\]
A row application loses \(\log(2D_{\rm m})\) in norm, hence
\(\log^2(2D_{\rm m})\) in squared energy.  This loss, the
controlled signed expansions, and the telescoping counts fit
\(G_0\).  The factors \((C\log(K+2))^{O(m)}\) and the binary
length choices also fit \(G_0\) by
\eqref{eq:h-budget-hierarchy}.  Thus no fixed power of \(K\)
enters \eqref{eq:h-stacks}.  Its powers in the component and
top bounds stay in the allowed \(P^C\) budgets.  The larger
\(D_*\) is used in the mode construction, Bessel, and counting arguments;
the final row adjoints in the two new lemmas use \(D_{\rm m}\).

\subsubsection{The fixed-input comparison on a short block}

For \(2<q<3\), put \(M_qg=(M(|g|^q))^{1/q}\), where \(M\) is
the uncentered Hardy--Littlewood maximal operator.  For every
integer \(N\ge1\), let \(C_N(q)\) be the least constant such that
\begin{equation}
 \sum_{I\in\mathcal A}|I||H_I(g_0,g_1,g_2,g_3)|
 \le C_N(q)\int_{\mathbb R}\prod_{j=0}^3M_qg_j
 \label{eq:h-fixed-local}
\end{equation}
for every dyadic lattice, finite collection \(\mathcal A\) in
at most \(N\) consecutive depths, admissible kernel family, and
four compactly supported scalar \(L^q\) functions \(g_j\) fixed
across the collection.  The \emph{Local absolute estimate} Lemma of H gives
\(C_N(q)\le CN\), so this constant is finite before its uniform
bound is used.

For our application choose \(N=d+1\), where \(d\) is the largest
original depth in Theorem~\ref{thm:depth-local}.  The summation
cells in each short-block piece are retained original active cells,
with their original depths in \(\{0,\ldots,d\}\).  Restarts only
change their assigned attempt.  Hull ancestors and completion
leaves, including a boundary level one step below a block, supply
the geometry and leaf norm conditions; they are not additional
form summation cells.

We use the following specialization of the
\emph{Short-block comparison} Lemma of \cite[Section~7]{H}.
Let a finite tree piece have internal summation cells in at most
\(K\) consecutive depths of this original \(N\)-depth range and
leaves in that block or one depth below.  Let deterministic
Hilbert-valued stacks \(\mathbf z_j\), supported on its root \(Q\),
satisfy \(\|\mathbf z_j\|_{2,L}\le s_j\) on every leaf \(L\).
Suppose scalar randomizations \(z_j^\sigma\) satisfy, for every
finite real \(a\ge1\),
\[
 \bigl(\mathbb E_\sigma|z_j^\sigma(y)|^a\bigr)^{1/a}
       \le C_a\|\mathbf z_j(y)\|
 \quad\text{for almost every }y,
\]
with \(C_a\) independent of the tree and \(K\).  Then, for
arbitrary measurable signs \(|\epsilon_I(\sigma)|\le1\) and any
subcollection of the internal cells,
\begin{equation}
 \mathbb E_\sigma\left|
   \sum_I|I|\epsilon_I(\sigma)
          H_I(z_0^\sigma,\ldots,z_3^\sigma)\right|
 \le (1+C_N(q))C_q(1+\log(K+2))^C
       (K+2)^{C_{19}(q-2)}|Q|\prod_j s_j.
 \label{eq:h-short-block}
\end{equation}
The degree \(C_{19}\) is absolute.  Products of at most two
independent Rademacher randomizations of a finite vector stack
satisfy the stated moment hypothesis.

Apply this comparison to the stacks in \eqref{eq:h-stacks},
the frozen complementary band stack, and the other frozen
complement.  The scalar inputs are sums of their physical
components.  Use independent signs \(\varepsilon_i,\eta_j\)
for the two position indices and signs \(\zeta_p\) for the
band index.  At the fixed offset \(\ell=q'-p\), index the complementary
coordinate \(q'=p+\ell\) by \(p\).  Randomize the banded side
by \(\varepsilon_i\zeta_p\), the untouched side by \(\eta_j\),
and that complementary band coordinate by \(\zeta_p\).  At a
use put \(\varepsilon_{i(s)}\eta_{j(s)}\) in the kernel.
Expectation selects the two positions and all matched bands,
leaving \((p-q')/P=O(P^{-1})\).  Iterated Khintchine
inequalities bound every finite scalar moment by the augmented
stack norm.

After fixing the signs, these randomized physical stack sums are
fixed scalar \(L^2\) functions.  As a randomized family, they satisfy
the pointwise moment bounds for every finite \(a\ge1\), controlled by
the deterministic stacks with the stated leaf sizes; the randomizations are not asserted to
lie in \(L^q\).  In H's \emph{Return to the continuous form by
bursts} step, fix the auxiliary random and fast parameters outside
a null set.  The burst inputs are then finite scalar linear
combinations of indicators of bounded intervals with finite
coefficients.  They are therefore compactly supported \(L^q\)
functions fixed across the original summation cells.  H applies
\eqref{eq:h-fixed-local} to these inputs with the original admissible
kernels and then integrates the auxiliary parameters.  The clipping
and mixed-maximal calculation in that proof gives the exact block
power \((K+2)^{11(q-2)/(q-1)}\le(K+2)^{11(q-2)}\); the remaining
block factors are logarithmic.

For the whole-hull calculation, let \(\Omega_\nu\) be the union
of the use cells for one full index \(\nu=(G,b,\chi)\).  Its
maximal use cells are disjoint roots, and the hull contains every cell
between those roots and their uses.  For the already fixed stack
define
\[
 m_\nu^{\rm full}(x)=
 \sup_{\substack{I\ {\rm in\ the\ full\ hull}\\x\in I}}
       \fint_I\|\mathrm{stack}_\nu\|^2.
\]
Set \(m_\nu^{\rm full}=0\) outside its hull.  The supremum is
defined on every whole hull cell, including branches below later cuts.
A zero-energy family contributes
nothing.  Otherwise suppose a family has total stack energy
\(E|R|\) with \(E>0\), and
\(\sum_\nu|\Omega_\nu|\le M_0|R|\), where
\(M_0\le\exp(C_qk^C)\) follows from the active cell/group path
counts and the auxiliary-choice overhead, with all groups and
blocks included in this support sum.  The dyadic maximal calculation
in
\cite[Section~7]{H} gives
\[
 \begin{aligned}
 \sum_\nu|\{m_\nu^{\rm full}>T\}|
 &\le\min(M_0|R|,CE|R|/T),\\
 \sum_\nu\int_R\min(m_\nu^{\rm full},C_1E)
 &\le CE(1+\log(2+C_1M_0))|R|.
 \end{aligned}
\]
For a fixed sufficiently large \(C_1\), first-crossing cuts on the
individual and then joint partial ancestor suprema yield the retained caps
\[
 \sum_\nu m_{a,\nu}(x)\le PG_0,\qquad
 \sum_\nu m_{\mathrm{untouched},\nu}(x)\le G_0.
\]
Here \(m_{a,\nu},m_{\mathrm{untouched},\nu}\) use the retained hull cells, still
taken whole.  Include every potentially needed hull ancestor
within the block before these common cuts.  At a crossing the
cell is removed before its value enters a retained cap; below a
cut only earlier retained ancestors contribute.  The caps
therefore hold also on portions of retained hull cells below a
cut, with the small relative length cost supplied by the full
maximal bounds.

On each retained hull, make the source sparse decomposition:
from a root \(Q\), stop where either squared side average first
exceeds sixteen times its average on \(Q\).  A zero root average gives
zero contribution.  The stopped cells
have total length at most \(|Q|/8\), so the disjoint remainders
\(E_Q\) have measure at least \(7|Q|/8\).  Complete each piece
by children immediately below its included cells.  A child's
squared average is at most twice its parent's, which gives the
leaf bounds in \eqref{eq:h-short-block}; the frozen complementary
sizes are bounded on the whole hull and its boundary leaves.
If \(a_Q,b_Q\) are the side sizes at the piece roots, the
whole-hull convention puts \(Q\) in both suprema at every
\(x\in E_Q\), even below a cut.  Hence
\[
 \begin{split}
 \sum_{\nu,Q}|Q|a_Qb_Q
 &\le C\int_R\sum_\nu\sqrt{m_{a,\nu}m_{\mathrm{untouched},\nu}}\\
 &\le C\int_R
       \left(\sum_\nu m_{a,\nu}\right)^{1/2}
       \left(\sum_\nu m_{\mathrm{untouched},\nu}\right)^{1/2}
 \le P^{1/2}G_0|R|.
 \end{split}
\]
This is the promised root-size cost from the stack energies and
the whole-hull cuts.  Combining it with
\eqref{eq:h-short-block} and the matched coefficient \(O(P^{-1})\)
gives the near-band form of Equation~\Heq{19} of
\cite[Section~7]{H}:
\begin{equation}
 (1+C_N(q))C_qP^{-1/2}G_0
                  (K+2)^{C_{19}(q-2)}|R|.
 \label{eq:h-near-band}
\end{equation}
It bounds the absolute sum of original cell contributions with
the prescribed band sum retained inside each contribution at a
fixed offset.  The logarithmic block factors fit \(G_0\); the
displayed power of \(K+2\) comes from the fixed-input
short-block comparison.

\subsection{The core calculation}

The remaining terms have at least three core positions.  We
specialize the conditional calculation in
\emph{Counts for the tested core tuples} and
\emph{Powered estimate for the core terms} of
\cite[Section~9]{H}.  Its good-height predicates, shifted
annular tests, pure-shift gap, and numerical separation and
height reserves remain exactly those of the source construction.
We check the inputs that can be affected by the changed endpoint
operands.

Fix a retained path point \(x_*\), a group, an upper height,
role types, frequency bins \((\mu,\nu,\Lambda)\), and amplitude
bins \(\sigma_j\).  Let \(N_s^G\) be H's exact count of tested
core representative-label tuples at depth \(s\) for these data.
In a core role put \(M_j^G=|\mathcal A_j^G|\).  In the optional
side role put
\[
 M_l^G=\sigma_l^{-2}\sum_{s:\,G\ {\rm used}}
                 \|S_{l,G,s}\|_{2,I_s}^2,\qquad
 M'=\max_jM_j^G.
\]
Omit groups with no nonzero count, so \(M'\ge1\).  H's cell
calculation gives, apart from total absolute error
\(O(e^{-k}|R|)\), terms of the form
\[
 C_L\mu\nu^{-L}\Lambda^{-L}
            \left(\prod_{j=0}^3\sigma_j\right)N_s^G.
\]
Here \(0<\mu\le1\le\nu\) and \(\Lambda\ge1\) are dyadic, with
\(\Lambda=1\) for four core positions.  The good-height masks
select the full core subtotals within the minimal CURRENT
components, with the coefficient floor already included in
\(\sigma_{\mathcal F}\).

The preceding arguments provide the three inputs to the source count
that depend on endpoint data.  First, a
representative label has the fixed value \(\omega_a(x_*)\).
Distinct labels in one role obey
\eqref{eq:h-representative-separation}.  Together with H's exact
tested-tuple constraints, this ensures that two labels in distinct
core roles determine at most one tested tuple, and
the source height-bin estimate bounds its scale multiplicity.
Second, the frozen-row, component-size, and top argument gives
\eqref{eq:h-top-count}, including the \(D_*\) cap in each core
role.  Third, a side role uses the actual same-cell function
\(S_{l,G,s}\) in its Fourier windows.  Parseval gives its
incidence bound
\[
 w_s=\sigma_l^{-2}\|S_{l,G,s}\|_{2,I_s}^2,
\]
whose sum is \(M_l^G\) and whose path budget is supplied by
\eqref{eq:h-side-cap}.

The remaining hypotheses of the
\emph{Sparse mass of separated progression tests} Lemma of
\cite[Section~9]{H} are exactly H's shifted frequency tests
and pure-shift gap, retained from its construction.  As in its proof of
\emph{Counts for the tested core tuples}, every successful
component contains a nonexceptional primary, whose
pure-horizontal alternatives exclude the required intermediate
range for every permitted shift.  This is a property of the
persistent mode graph and its parameters.  The finite Lemma
therefore applies to the persistent labels and the actual
side incidences and gives, for an absolute \(p_0\in(0,1)\),
\begin{equation}
 \begin{gathered}
 \sum_sN_s^G\le L_0\min_{i\ne j}M_i^GM_j^G,\qquad
 L_0\le P^C(1+\log(\nu/\mu))^C,\\
 \sum_s(N_s^G)^{p_0}
 \le P^C(1+\nu+\mu^{-1}+\Lambda)^C(M')^{2p_0}.
 \end{gathered}
 \label{eq:h-core-counts}
\end{equation}
This source calculation identifies representative labels
across depths and uses same-cell side data; coefficients frozen
in different blocks may be unrelated.

Finally \eqref{eq:h-top-count} and \eqref{eq:h-side-cap} give
\(\sum_GM_j^G\le P^C\sigma_j^{-2}\) in every role, with the
additional \(D_*\) cap in a core role.  H interpolates the two
inequalities in \eqref{eq:h-core-counts}, keeps \(N_s^G\)
inside its power, and sums the amplitude and frequency bins.
Its \emph{Powered estimate for the core terms},
Equation~\Heq{37} of \cite[Section~9]{H}, gives absolute
constants \(\varepsilon>0\) and \(C_4\), with
\(p=1-\varepsilon>2/3\), and total contribution
\begin{equation}
 C_q\eta^\varepsilon P^{C_4}|R|+O(e^{-k}|R|)
 \label{eq:h-core-estimate}
\end{equation}
for the terms with at least three core positions on this
attempt's retained cells.  The original cell minimum is applied
before subadditivity and the group count remains inside the
power.  Numerical errors are estimated absolutely.  The
conditional source calculation thus uses the supplied persistent
labels, full core-subtotal masks, role counts, and actual side
Fourier data, with no further cross-depth estimate on an
original input.
